% GENERATED FILE. Edit canonical lessons, then run npm run generate:books.
% canonical-source-hash: fnv1a64:74da2e7463f9355b
% canonical-lessons: SA-C349-nyayalayah, SA-C349-prativadi, SA-C349-aparadhyati, SA-C349-nigrhnati, SA-C349-sapathah

\chapter{In Court}
\label{ch:sa-in-court}
\hlchaptermodality{349}

\begin{chapteropening}
\textbf{By the end of this chapter:} \emph{I can follow what happens to a defendant in court.}

Complete the last lesson of chapter 349: I can follow what happens to a defendant in court.
\end{chapteropening}

\section[\texorpdfstring{\sk{न्यायालयः} (nyāyālayaḥ)}{न्यायालयः (nyāyālayaḥ)}]{\sk{न्यायालयः} (nyāyālayaḥ) — a court of law}

\label{lesson:SA-C349-nyayalayah}



\begin{quote}
Before the new one: say the Sanskrit for partiality, then the Sanskrit for breaks (a rule).
\end{quote}

\subsection*{You'll want to know: \sk{न्यायालयः}}
\textbf{\sk{न्यायालयः}} — \emph{nyāyālayaḥ} — \textquotedblleft{}a court of law\textquotedblright{}.

\begin{quote}
\sk{अधिवक्ता} \sk{प्रातः} \sk{न्यायालयं} \sk{गच्छति।}
\end{quote}

\begin{quote}
\emph{adhivaktā prātaḥ nyāyālayaṃ gacchati.}
\end{quote}

\begin{quote}
The lawyer goes to court in the morning.
\end{quote}

A masculine noun.

\begin{cousinweb}
\textbf{\sk{न्यायालयः}} (\emph{nyāyālayaḥ}) is a modern word built from \emph{nyāya}, \textquotedblleft{}rule, right manner\textquotedblright{}, and \emph{ālaya}, \textquotedblleft{}a house, an abode\textquotedblright{}; MW records the parts but not the compound (MW).
\end{cousinweb}

\begin{grammarlens}[title={What you've built}]
One more word for giving reasons and opinions.
\end{grammarlens}

\subsection*{Your turn}
\begin{itemize}
\raggedright
  \item \emph{Say it:} \emph{nyāyālayaḥ}
  \item \emph{Say it:} \emph{nyāyālayaḥ}, once more
  \item \emph{Say it:} \emph{adhivaktā prātaḥ nyāyālayaṃ gacchati.} — The lawyer goes to court in the morning.
  \item \emph{Recall:} say the Sanskrit for partiality, then the Sanskrit for breaks (a rule), then say \emph{nyāyālayaḥ} again
\end{itemize}

\begin{tcolorbox}[breakable,colback=teal!4,colframe=teal!35!black,title={Before you move on}]
What does \sk{न्यायालयः} mean? (\textquotedblleft{}A court of law\textquotedblright{}.) Say it once more.
\end{tcolorbox}
\section[\texorpdfstring{\sk{प्रतिवादी} (prativādī)}{प्रतिवादी (prativādī)}]{\sk{प्रतिवादी} (prativādī) — a defendant}

\label{lesson:SA-C349-prativadi}



\begin{quote}
Before the new one: say the Sanskrit for breaks (a rule), then the Sanskrit for a court of law.
\end{quote}

\subsection*{You'll want to know: \sk{प्रतिवादी}}
\textbf{\sk{प्रतिवादी}} — \emph{prativādī} — \textquotedblleft{}a defendant\textquotedblright{}.

\begin{quote}
\sk{प्रतिवादी} \sk{न्यायाधीशाय} \sk{सत्यं} \sk{वदति।}
\end{quote}

\begin{quote}
\emph{prativādī nyāyādhīśāya satyaṃ vadati.}
\end{quote}

\begin{quote}
The defendant tells the judge the truth.
\end{quote}

A masculine or feminine noun, by the person meant.

\begin{cousinweb}
\textbf{\sk{प्रतिवादी}} (\emph{prativādī}) is built from \emph{prati}, \textquotedblleft{}against\textquotedblright{}, and \emph{vādin}, \textquotedblleft{}speaking\textquotedblright{}, from the root √vad; MW gives \emph{prati-vādin} \textquotedblleft{}an opponent\textquotedblright{} and \textquotedblleft{}a defendant, a respondent\textquotedblright{} (MW).
\end{cousinweb}

\begin{grammarlens}[title={What you've built}]
One more word for giving reasons and opinions.
\end{grammarlens}

\subsection*{Your turn}
\begin{itemize}
\raggedright
  \item \emph{Say it:} \emph{prativādī}
  \item \emph{Say it:} \emph{prativādī}, once more
  \item \emph{Say it:} \emph{prativādī nyāyādhīśāya satyaṃ vadati.} — The defendant tells the judge the truth.
  \item \emph{Recall:} say the Sanskrit for breaks (a rule), then the Sanskrit for a court of law, then say \emph{prativādī} again
\end{itemize}

\begin{tcolorbox}[breakable,colback=teal!4,colframe=teal!35!black,title={Before you move on}]
What does \sk{प्रतिवादी} mean? (\textquotedblleft{}A defendant\textquotedblright{}.) Say it once more.
\end{tcolorbox}
\section[\texorpdfstring{\sk{अपराध्यति} (aparādhyati)}{अपराध्यति (aparādhyati)}]{\sk{अपराध्यति} (aparādhyati) — does wrong; offends}

\label{lesson:SA-C349-aparadhyati}



\begin{quote}
Before the new one: say the Sanskrit for a court of law, then the Sanskrit for a defendant.
\end{quote}

\subsection*{You'll want to know: \sk{अपराध्यति}}
\textbf{\sk{अपराध्यति}} — \emph{aparādhyati} — \textquotedblleft{}does wrong; offends\textquotedblright{}.

\begin{quote}
\sk{यः} \sk{अपराध्यति}, \sk{सः} \sk{दण्डं} \sk{प्राप्नोति।}
\end{quote}

\begin{quote}
\emph{yaḥ aparādhyati, saḥ daṇḍaṃ prāpnoti.}
\end{quote}

\begin{quote}
Whoever does wrong is punished.
\end{quote}

Past (he …): \sk{सः} \sk{अपराद्धवान्}, \emph{saḥ aparāddhavān}.

\begin{cousinweb}
\textbf{\sk{अपराध्यति}} (\emph{aparādhyati}) is built on \emph{apa}, \textquotedblleft{}away\textquotedblright{}, and the root √rādh, \textquotedblleft{}to succeed\textquotedblright{}; MW gives \emph{apa-√rādh} \textquotedblleft{}to wrong, to offend against, to sin\textquotedblright{} (MW).
\end{cousinweb}

\begin{grammarlens}[title={What you've built}]
One more word for giving reasons and opinions.
\end{grammarlens}

\subsection*{Your turn}
\begin{itemize}
\raggedright
  \item \emph{Say it:} \emph{aparādhyati}
  \item \emph{Say it:} \emph{aparādhyati}, once more
  \item \emph{Say it:} \emph{yaḥ aparādhyati, saḥ daṇḍaṃ prāpnoti.} — Whoever does wrong is punished.
  \item \emph{Recall:} say the Sanskrit for a court of law, then the Sanskrit for a defendant, then say \emph{aparādhyati} again
\end{itemize}

\begin{tcolorbox}[breakable,colback=teal!4,colframe=teal!35!black,title={Before you move on}]
What does \sk{अपराध्यति} mean? (\textquotedblleft{}Does wrong; offends\textquotedblright{}.) Say it once more.
\end{tcolorbox}
\section[\texorpdfstring{\sk{निगृह्णाति} (nigṛhṇāti)}{निगृह्णाति (nigṛhṇāti)}]{\sk{निगृह्णाति} (nigṛhṇāti) — restrains; punishes}

\label{lesson:SA-C349-nigrhnati}



\begin{quote}
Before the new one: say the Sanskrit for a defendant, then the Sanskrit for does wrong.
\end{quote}

\subsection*{You'll want to know: \sk{निगृह्णाति}}
\textbf{\sk{निगृह्णाति}} — \emph{nigṛhṇāti} — \textquotedblleft{}restrains; punishes\textquotedblright{}.

\begin{quote}
\sk{राजा} \sk{चौरान्} \sk{निगृह्णाति।}
\end{quote}

\begin{quote}
\emph{rājā caurān nigṛhṇāti.}
\end{quote}

\begin{quote}
The king restrains thieves.
\end{quote}

Past (he …): \sk{सः} \sk{चौरं} \sk{निगृहीतवान्}, \emph{saḥ cauraṃ nigṛhītavān}.

\begin{cousinweb}
\textbf{\sk{निगृह्णाति}} (\emph{nigṛhṇāti}) is built on \emph{ni}, \textquotedblleft{}down\textquotedblright{}, and the root √grah, \textquotedblleft{}to seize\textquotedblright{}; MW gives \emph{ni-√grah} \textquotedblleft{}to hold back, restrain, curb, punish\textquotedblright{} (MW).
\end{cousinweb}

\begin{grammarlens}[title={What you've built}]
One more word for giving reasons and opinions.
\end{grammarlens}

\subsection*{Your turn}
\begin{itemize}
\raggedright
  \item \emph{Say it:} \emph{nigṛhṇāti}
  \item \emph{Say it:} \emph{nigṛhṇāti}, once more
  \item \emph{Say it:} \emph{rājā caurān nigṛhṇāti.} — The king restrains thieves.
  \item \emph{Recall:} say the Sanskrit for a defendant, then the Sanskrit for does wrong, then say \emph{nigṛhṇāti} again
\end{itemize}

\begin{tcolorbox}[breakable,colback=teal!4,colframe=teal!35!black,title={Before you move on}]
What does \sk{निगृह्णाति} mean? (\textquotedblleft{}Restrains; punishes\textquotedblright{}.) Say it once more.
\end{tcolorbox}
\section[\texorpdfstring{\sk{शपथः} (śapathaḥ)}{शपथः (śapathaḥ)}]{\sk{शपथः} (śapathaḥ) — an oath}

\label{lesson:SA-C349-sapathah}



\begin{quote}
Before the new one: say the Sanskrit for does wrong, then the Sanskrit for restrains.
\end{quote}

\subsection*{You'll want to know: \sk{शपथः}}
\textbf{\sk{शपथः}} — \emph{śapathaḥ} — \textquotedblleft{}an oath\textquotedblright{}.

\begin{quote}
\sk{सः} \sk{सत्यस्य} \sk{शपथं} \sk{करोति।}
\end{quote}

\begin{quote}
\emph{saḥ satyasya śapathaṃ karoti.}
\end{quote}

\begin{quote}
He swears an oath to tell the truth.
\end{quote}

A masculine noun.

\begin{cousinweb}
\textbf{\sk{शपथः}} (\emph{śapathaḥ}) comes from the root √śap, \textquotedblleft{}to curse, to swear\textquotedblright{}; MW gives \emph{śapatha} \textquotedblleft{}an oath, a vow\textquotedblright{} (MW).
\end{cousinweb}

\begin{grammarlens}[title={What you've built}]
One more word for giving reasons and opinions.
\end{grammarlens}

\subsection*{Your turn}
\begin{itemize}
\raggedright
  \item \emph{Say it:} \emph{śapathaḥ}
  \item \emph{Say it:} \emph{śapathaḥ}, once more
  \item \emph{Say it:} \emph{saḥ satyasya śapathaṃ karoti.} — He swears an oath to tell the truth.
  \item \emph{Recall:} say the Sanskrit for does wrong, then the Sanskrit for restrains, then say \emph{śapathaḥ} again
\end{itemize}

\begin{tcolorbox}[breakable,colback=teal!4,colframe=teal!35!black,title={Before you move on}]
What does \sk{शपथः} mean? (\textquotedblleft{}An oath\textquotedblright{}.) Say it once more.
\end{tcolorbox}
